\subsection{Pengurangan Fungsi}
\begin{definition}[Pengurangan Fungsi]
  \[ (f - g)(x) = f(x) - g(x) \]
\end{definition}

\begin{example}[Pengurangan Fungsi]
  \hfill 
  \begin{enumerate}
  \item Jika $f(x) = 5x + 7$ dan $g(x) = 3x - 2$:
    \[
    \begin{aligned}
      (f - g)(x) &= (5x + 7) - (3x - 2) \\
      &= 5x + 7 - 3x + 2 \\
      &= 2x + 9
    \end{aligned}
    \]
  \item Jika $f(x) = 2x^2 - 6x + 10$ dan $g(x) = x^2 + 4x - 5$:
    \[
    \begin{aligned}
      (f - g)(x) &= (2x^2 - 6x + 10) - (x^2 + 4x - 5) \\
      &= 2x^2 - 6x + 10 - x^2 - 4x + 5 \\
      &= x^2 - 10x + 15
    \end{aligned}
    \]
  \end{enumerate}
\end{example}

\subsection{Perkalian Fungsi}
\begin{definition}[Perkalian Fungsi]
  Konsep distributif:
  \[ a(b + c) = ab + ac \]
\end{definition}

\begin{example}[Perkalian Fungsi]
  \hfill \\
  Misalkan $f(x) = 2x$ dan $g(x) = x - 7$:
  \[
  \begin{aligned}
    (f \cdot g)(x) &= 2x \cdot (x - 7) \\
    &= 2x^2 - 14x
  \end{aligned}
  \]
  \[
  \begin{aligned}
    (g \cdot f)(x) &= (x - 7) \cdot 2x \\
    &= 2x^2 - 14x
  \end{aligned}
  \]
\end{example}

\begin{definition}[Perkalian Fungsi Dua Suku]
  \[ (a + b) \cdot (c + d) = ac + ad + bc + bd \]
  Khusus untuk kuadrat dua suku:
  \[ (a + b)^2 = a^2 + 2ab + b^2 \]
\end{definition}

\begin{example}[Perkalian Fungsi Dua Suku]
  \hfill \\
  Misalkan $f(x) = 2x - 3$ dan $g(x) = 3x + 4$:
  \[
  \begin{aligned}
    (f \cdot g)(x) &= (2x - 3) \cdot (3x + 4) \\
    &= 6x^2 + 8x - 9x - 12 \\
    &= 6x^2 - x - 12
  \end{aligned}
  \]
\end{example}

\textbf{Note:} Aturan saat mengalikan:
\begin{enumerate}
\item Gunakan sifat distributif.
\item Jika variabelnya sama, pangkat dijumlahkan ($x^m \cdot x^n = x^{m + n}$).
\item Kalikan angka dengan angka.
\item Kelompokkan suku sejenis.
\item Sederhanakan.
\end{enumerate}
